% 章首引语及所据原文出处。
\DeclareBookEpigraph{preface}%
  {智慧是经验的女儿。}%
  {列奥纳多·达·芬奇，《笔记》§1150}%
  {https://www.gutenberg.org/ebooks/5000}

\DeclareBookEpigraph{introduction}%
  {世界是事实的总和，而非事物的总和。}%
  {路德维希·维特根斯坦，《逻辑哲学论》1.1}%
  {https://www.gutenberg.org/files/5740/5740-pdf.pdf}

\DeclareBookEpigraph{overview}%
  {观念的次序和联系，与事物的次序和联系相同。}%
  {巴鲁赫·斯宾诺莎，《伦理学》第二部分，命题七}%
  {https://www.gutenberg.org/ebooks/3800}

\DeclareBookEpigraph{01}%
  {程序必须写给人阅读，机器执行只是附带的目的。}%
  {阿贝尔森、杰拉尔德·萨斯曼与朱莉·萨斯曼，《计算机程序的构造和解释》初版序}%
  {https://sarabander.github.io/sicp/html/Preface-1e.xhtml}

\DeclareBookEpigraph{02}%
  {一切关于事实的推理，似乎都建立在因果关系之上。}%
  {大卫·休谟，《人类理解研究》第四节，§22}%
  {https://www.gutenberg.org/ebooks/9662}

\DeclareBookEpigraph{03}%
  {名称指谓对象，对象就是名称的意义。}%
  {路德维希·维特根斯坦，《逻辑哲学论》3.203}%
  {https://www.gutenberg.org/files/5740/5740-pdf.pdf}

\DeclareBookEpigraph{04}%
  {在分布式系统中，要认识到事件发生的次序只是一个偏序。}%
  {莱斯利·兰伯特，《时间、时钟与分布式系统中的事件排序》结论}%
  {https://lamport.azurewebsites.net/pubs/time-clocks.pdf}

\DeclareBookEpigraph{05}%
  {我们为自己构造事实的图像。}%
  {路德维希·维特根斯坦，《逻辑哲学论》2.1}%
  {https://www.gutenberg.org/files/5740/5740-pdf.pdf}

\DeclareBookEpigraph{06}%
  {程序测试可以用来显示错误的存在，却永远不能显示错误的不存在。}%
  {艾兹赫尔·戴克斯特拉，《结构化程序设计札记》EWD249}%
  {https://www.cs.utexas.edu/~EWD/transcriptions/EWD02xx/EWD249/EWD249.html}

\DeclareBookEpigraph{part1}%
  {从最简单、最容易认识的对象开始，一点一点、逐步上升，直到认识较为复杂的对象。}%
  {勒内·笛卡尔，《方法谈》第二部分，第三条准则}%
  {https://www.gutenberg.org/ebooks/59}

\DeclareBookEpigraph{07}%
  {把任务分开是好事，同时，我们还得把分开后的端头重新接起来。}%
  {艾兹赫尔·戴克斯特拉，《结构化程序设计札记》EWD249}%
  {https://www.cs.utexas.edu/~EWD/transcriptions/EWD02xx/EWD249/EWD249.html}

\DeclareBookEpigraph{08}%
  {每个物体都保持静止或匀速直线运动的状态，除非外加的力迫使它改变这种状态。}%
  {艾萨克·牛顿，《自然哲学的数学原理》运动第一定律}%
  {https://www.gutenberg.org/ebooks/28233}

\DeclareBookEpigraph{09}%
  {通信的根本问题，是在一个地点精确或近似地重现另一个地点选定的消息。}%
  {克劳德·香农，《通信的数学理论》引言}%
  {https://people.math.harvard.edu/~ctm/home/text/others/shannon/entropy/entropy.pdf}

\DeclareBookEpigraph{10}%
  {我们通过确立约定的接口来控制复杂性，使标准而充分理解的部件能够自由组合，构成系统。}%
  {阿贝尔森、杰拉尔德·萨斯曼与朱莉·萨斯曼，《计算机程序的构造和解释》初版序}%
  {https://sarabander.github.io/sicp/html/Preface-1e.xhtml}

\DeclareBookEpigraph{part2}%
  {整体大于部分。}%
  {欧几里得，《几何原本》第一卷，公理五}%
  {https://mathcs.clarku.edu/~djoyce/elements/bookI/cn.html}

\DeclareBookEpigraph{11}%
  {现在孕育着未来。}%
  {戈特弗里德·威廉·莱布尼茨，《单子论》§22}%
  {https://www.gutenberg.org/ebooks/17641}

\DeclareBookEpigraph{12}%
  {好的记号使头脑摆脱一切不必要的工作，能够集中精力处理更高深的问题。}%
  {阿尔弗雷德·诺思·怀特海，《数学导论》第五章}%
  {https://www.gutenberg.org/files/41568/41568-pdf.pdf}

\DeclareBookEpigraph{13}%
  {数学是一门给不同事物以相同名称的艺术。这些事物虽在内容上不同，在形式上却应当相似。}%
  {亨利·庞加莱，《科学与方法》第一编第二章}%
  {https://www.gutenberg.org/ebooks/39713}

\DeclareBookEpigraph{14}%
  {我们应当把宇宙的现在状态，看作先前状态的结果，以及随后状态的原因。}%
  {皮埃尔-西蒙·拉普拉斯，《概率的哲学论》第二章}%
  {https://www.gutenberg.org/ebooks/58881}

\DeclareBookEpigraph{15}%
  {因此，明智的人会使自己的信念与证据相称。}%
  {大卫·休谟，《人类理解研究》第十节，§87}%
  {https://www.gutenberg.org/ebooks/9662}

\DeclareBookEpigraph{part3}%
  {真理是时间唯一的女儿。}%
  {列奥纳多·达·芬奇，《笔记》§1152}%
  {https://www.gutenberg.org/ebooks/5000}

\DeclareBookEpigraph{16}%
  {把所考察的每一个困难，尽可能分成充分解决它所需要的若干部分。}%
  {勒内·笛卡尔，《方法谈》第二部分，第二条准则}%
  {https://www.gutenberg.org/ebooks/59}

\DeclareBookEpigraph{17}%
  {我们通过构造抽象，在适当的时候隐藏细节，以此控制复杂性。}%
  {阿贝尔森、杰拉尔德·萨斯曼与朱莉·萨斯曼，《计算机程序的构造和解释》初版序}%
  {https://sarabander.github.io/sicp/html/Preface-1e.xhtml}

\DeclareBookEpigraph{18}%
  {形式是这种内在意义的外在表达。}%
  {瓦西里·康定斯基，《论艺术的精神》第六章}%
  {https://www.gutenberg.org/ebooks/5321}

\DeclareBookEpigraph{19}%
  {如同在数学中一样，在自然哲学中，用分析方法探究困难问题，应当先于综合方法。}%
  {艾萨克·牛顿，《光学》第三卷，疑问三十一}%
  {https://www.gutenberg.org/ebooks/33504}

\DeclareBookEpigraph{20}%
  {文明的进步，在于扩大那些我们无须思索便能完成的重要操作的数量。}%
  {阿尔弗雷德·诺思·怀特海，《数学导论》第五章}%
  {https://www.gutenberg.org/files/41568/41568-pdf.pdf}

\DeclareBookEpigraph{21}%
  {我们只能看见前方不远的地方，但已能看见那里有许多工作要做。}%
  {艾伦·图灵，《计算机器与智能》结尾}%
  {https://www.csee.umbc.edu/courses/471/papers/turing.pdf}

\DeclareBookEpigraph{part4}%
  {每一件艺术作品都是它所属时代的孩子，而且在许多情况下，也是我们情感的母亲。}%
  {瓦西里·康定斯基，《论艺术的精神》第一章}%
  {https://www.gutenberg.org/ebooks/5321}

\DeclareBookEpigraph{closing}%
  {科学由事实构成，如同房屋由石头建成；一堆事实还构不成科学，正如一堆石头还构不成房屋。}%
  {亨利·庞加莱，《科学与假设》第九章}%
  {https://www.gutenberg.org/ebooks/39713}
